		\subsection{\texttt{segtree}：单点改 + 区间查（非递归）}

		维护满足\textbf{结合律}的二元运算 \texttt{op} 与其\textbf{单位元} \texttt{e}（即幺半群 monoid）上的序列，支持单点赋值 $O(\log n)$、区间求积 $O(\log n)$、以及在线段树上二分 $O(\log n)$。

		\textbf{接口一览}（$n$ 为序列长度，下标 0-based）：

		\begin{itemize}
			\item \texttt{segtree<S, op, e> seg(n)}：长度 $n$，全部初始化为 \texttt{e()}。
			\item \texttt{segtree<S, op, e> seg(v)}：由 \texttt{vector<S> v} 建树，$O(n)$。
			\item \texttt{seg.set(p, x)}：令 $a_p \leftarrow x$。
			\item \texttt{seg.get(p)}：返回 $a_p$。
			\item \texttt{seg.prod(l, r)}：返回 $op(a_l, \dots, a_{r-1})$；$l = r$ 时返回 \texttt{e()}。
			\item \texttt{seg.all\_prod()}：返回全区间积，$O(1)$。
			\item \texttt{seg.max\_right(l, f)}：返回最大的 $r$，使 $f(op(a_l, \dots, a_{r-1}))$ 为真。要求 \texttt{f} 单调（真在前假在后）且 \texttt{f(e())} 为真。
			\item \texttt{seg.min\_left(r, f)}：返回最小的 $l$，使 $f(op(a_l, \dots, a_{r-1}))$ 为真。同样要求 \texttt{f} 单调且 \texttt{f(e())} 为真。
		\end{itemize}

		\textbf{三个易错点}：

		\begin{enumerate}
			\item \texttt{op} 与 \texttt{e} 必须是\textbf{自由函数}（或函数指针 / 无捕获 lambda 转成的函数指针），\textbf{不能}是带捕获的 lambda——它们是模板参数，编译期就要定死。
			\item \texttt{e()} 必须是真单位元：求最大值时用足够小的哨兵（如 \texttt{-4e18}）而非 $0$，否则负数序列查询结果错。
			\item \texttt{prod} 是 $[l, r)$ 左闭右开；从 1-based 闭区间 $[L, R]$ 换算过来是 \texttt{prod(L - 1, R)}。
		\end{enumerate}

		\textbf{示例一：区间最大值}（单点赋值 + 区间查最大，完整可提交骨架）

		\begin{minted}{cpp}
/* 题意: n 个数, q 次操作
 *   1 p v -> 把第 p 个数改成 v  (1-based)
 *   2 l r -> 输出 [l, r] 里的最大值 (闭区间)
 */
using S = ll;   // 每个结点存什么: 这里就是值本身

// op: 两段的答案怎么合并成一段的答案
// 必须满足结合律 —— max 满足:
//   max(max(a,b),c) == max(a,max(b,c))
S op_max(S a, S b) { return max(a, b); }

// e: op 的单位元, 要求 op(e(), x) == x 恒成立
// 所以取「不可能出现的极小值」当哨兵。
// 千万别写 0: 全是负数时查询会错答成 0
S e_max() { return -4e18; }

int main() {
    int n, q;
    cin >> n >> q;

    vector<S> a(n);            // 全程 0-based
    for (auto &x: a) cin >> x;

    // 用 vector 一次性建树是 O(n);
    // 先 segtree(n) 再逐个 set 是 O(n log n),
    // 能一次建好就不要循环 set
    segtree<S, op_max, e_max> seg(a);

    while (q--) {
        int t;
        cin >> t;
        if (t == 1) {          // 单点赋值
            int p;
            ll v;
            cin >> p >> v;
            seg.set(p - 1, v); // 题面 1-based, 减 1
        } else {               // 区间查最大
            int l, r;
            cin >> l >> r;
            // prod 是左闭右开 [l, r), 而题面给的
            // 是闭区间 [l, r] -> 写成 prod(l-1, r)
            cout << seg.prod(l - 1, r) << "\n";
        }
    }
}
		\end{minted}

		\textbf{示例二：区间和 + 线段树上二分}。求「从 $l$ 出发最长的前缀，使区间和不超过 $x$」——这是 \texttt{max\_right} 的典型用途，把「外层二分 + 内层查询」的 $O(\log^2 n)$ 压到 $O(\log n)$：

		\begin{minted}{cpp}
using S = ll;
// 求和版三件套: 加法结合, 单位元是 0
S op_sum(S a, S b) { return a + b; }
S e_sum() { return 0; }

int main() {
    int n;
    ll x;
    cin >> n >> x;
    vector<S> a(n);
    for (auto &v: a) cin >> v;
    segtree<S, op_sum, e_sum> seg(a);

    int l = 0;   // 从下标 l 出发

    // 谓词 f 可以是「带捕获的 lambda」——
    // 它是普通函数参数, 不是模板参数,
    // 这点跟 op / e 的限制不同。
    // 要求 f 单调: 真在前、假在后, 且 f(e()) 真。
    // 这里 sum <= x 单调的前提是元素非负
    // (有负数时前缀和会上下摆, 谓词不单调,
    //  max_right 的返回值就没有意义了)
    int r = seg.max_right(l, [&](S sum) {
        return sum <= x;
    });

    // r 是「最大的下标」, 使
    //   a[l] + a[l+1] + ... + a[r-1] <= x
    // 所以能取的元素个数 = r - l
    cout << r - l << "\n";

    // 对称接口: min_left(r, f) 从右往左找
    // 最小的 l 使 a[l..r-1] 满足 f
    int l2 = seg.min_left(n, [&](S sum) {
        return sum <= x;
    });
    cout << n - l2 << "\n";  // 后缀能取几个
}
		\end{minted}

		\textbf{模板本体}：

		\begin{minted}{cpp}
/* AtCoder Library v1.6 atcoder/segtree 自包含移植
 * 改动: 去 #ifndef 守卫与 namespace atcoder,
 *      去 std:: 前缀, internal_bit 的
 *      bit_ceil / countr_zero 内联成私有 static
 * 约定: 下标 0-based, 区间 [l, r) 左闭右开
 */
template<class S, S (*op)(S, S), S (*e)()>
struct segtree {
public:
    segtree() : segtree(0) {}

    explicit segtree(int n)
        : segtree(vector<S>(n, e())) {}

    explicit segtree(const vector<S> &v)
        : _n(int(v.size())) {
        size = (int) bit_ceil((unsigned) _n);
        log = countr_zero((unsigned) size);
        d = vector<S>(2 * size, e());
        for (int i = 0; i < _n; i++)
            d[size + i] = v[i];
        for (int i = size - 1; i >= 1; i--)
            update(i);
    }

    // a[p] = x
    void set(int p, S x) {
        assert(0 <= p && p < _n);
        p += size;
        d[p] = x;
        for (int i = 1; i <= log; i++)
            update(p >> i);
    }

    // 返回 a[p]
    S get(int p) const {
        assert(0 <= p && p < _n);
        return d[p + size];
    }

    // 返回 op(a[l], ..., a[r - 1])
    S prod(int l, int r) const {
        assert(0 <= l && l <= r && r <= _n);
        S sml = e(), smr = e();
        l += size, r += size;
        while (l < r) {
            if (l & 1) sml = op(sml, d[l++]);
            if (r & 1) smr = op(d[--r], smr);
            l >>= 1, r >>= 1;
        }
        return op(sml, smr);
    }

    S all_prod() const { return d[1]; }

    // 最大的 r 使 f(op(a[l..r-1])) 为真
    template<class F>
    int max_right(int l, F f) const {
        assert(0 <= l && l <= _n);
        assert(f(e()));
        if (l == _n) return _n;
        l += size;
        S sm = e();
        do {
            while (l % 2 == 0) l >>= 1;
            if (!f(op(sm, d[l]))) {
                while (l < size) {
                    l = 2 * l;
                    if (f(op(sm, d[l]))) {
                        sm = op(sm, d[l]);
                        l++;
                    }
                }
                return l - size;
            }
            sm = op(sm, d[l]);
            l++;
        } while ((l & -l) != l);
        return _n;
    }

    // 最小的 l 使 f(op(a[l..r-1])) 为真
    template<class F>
    int min_left(int r, F f) const {
        assert(0 <= r && r <= _n);
        assert(f(e()));
        if (r == 0) return 0;
        r += size;
        S sm = e();
        do {
            r--;
            while (r > 1 && (r % 2)) r >>= 1;
            if (!f(op(d[r], sm))) {
                while (r < size) {
                    r = 2 * r + 1;
                    if (f(op(d[r], sm))) {
                        sm = op(d[r], sm);
                        r--;
                    }
                }
                return r + 1 - size;
            }
            sm = op(d[r], sm);
        } while ((r & -r) != r);
        return 0;
    }

private:
    int _n, size, log;
    vector<S> d;

    // 原 atcoder::internal::bit_ceil
    static unsigned bit_ceil(unsigned n) {
        unsigned x = 1;
        while (x < n) x *= 2;
        return x;
    }

    // 原 atcoder::internal::countr_zero
    static int countr_zero(unsigned n) {
        return __builtin_ctz(n);
    }

    void update(int k) {
        d[k] = op(d[2 * k], d[2 * k + 1]);
    }
};
		\end{minted}

		\textbf{与本模板其它线段树的取舍}：ACL 版是非递归实现，常数小、代码短，但\textbf{只支持单点修改}；要区间修改（区间加 / 区间赋值）请用下一小节的 \texttt{lazy\_segtree}，或主板子数据结构章的递归式懒标记线段树。ACL 版另一个优势是 \texttt{max\_right} / \texttt{min\_left} 这对树上二分接口，主板子的线段树模板没有等价物。
